% s-358-p358-m002; original PDF 358, 359, 360, 361, 362, 363
\Aphorism{s-358-p358-m002}{406}
\Sperr{Die Wiederkehr des Gleichen.} Wieviel man auf Nietzsches glänzende Speculation von der ewigen Wiederkunft geben will, wird dem Einzelnen zu überlassen sein; auf die Art Beweis, die hinterher dazu ersonnen und jetzt in seinem Nachlass zum Vorschein gekommen ist, rathe ich jedenfalls nichts zu geben. Beweise sind, bei allen Philosophen, das Verdächtige und überdies Langweilige, eine blosse Stubenhockerei aus missverstandener Gewissenhaftigkeit, ein gelehrt sein sollender Formalismus, den der Autor sich so wenig schenken zu dürfen glaubt, so gern ihn der Leser schenken würde. Geistreich, unterhaltend, bisweilen grandios sind höchstens die philosophischen Gedanken, die aufblitzenden und schon wieder verschwindenden Perspectiven, jene Plötzlichkeiten und Abgrundsblicke, die uns mitten im Nebel etwas vom Hochgebirge verrathen; aber der Nebel — die sogenannten Beweise — taugt nichts. Leider hat sich die Philosophie, namentlich die deutsche, immer zu sehr in der Nähe der Wissenschaft aufgehalten, statt sich unbefangen als das zu geben, was sie ist, als Religion oder Kunst auf Umwegen; daher die Nachahmung gelehrter Manieren, das Beweisenwollen, wo nichts zu beweisen ist, sei es \Foreign{more geometrico} oder aus Kategorientafeln oder „nach inductiver Methode“ oder mit Hülfe \Sperr{atomistischer} Betrachtungen. Soviel sich nämlich aus dem Entwurf „Die Wiederkehr des Gleichen“ ersehen lässt, scheint Nietzsche der Unendlichkeit der Zeit eine materialistisch gefasste Endlichkeit des Zeitinhalts gegenüberzustellen und daraus die nothwendige Wiederholung gleicher Zeitstrecken zu folgern. Wenn Raum, Kraft, Materie, Anzahl der Atome und ihrer möglichen Gruppirungen endliche Grössen sind, die Zeit aber eine unendliche, so ist in der That der Weltinhalt zu arm, um durch einmaliges Abrollen die Zeit auszufüllen, und der Dirigent müsste, um weiter zu spielen, das Stück ewig repetiren. An diesem Schluss ist weiter nichts hinfällig als die Voraussetzung, dass der mögliche Zeitinhalt gegenüber der unendlichen Zeit eine endliche Grösse sei; er ist in Wirklichkeit ebenfalls eine unendliche, und zwar eine viel umfassendere, umfänglichere, eine Unendlichkeit \Sperr{höherer Dimension}. Ohne diese der Mathematik entlehnte Unterscheidung der Mannigfaltigkeiten verschiedener Dimension, der Unendlichkeiten verschiedenen Rangs und Umfangs sollte sich Niemand an erkenntnisstheoretische Dinge heranwagen; wie will man mit der Unendlichkeit umgehen, wenn man in ihr bloss summarisch den Gegensatz des Endlichen, nicht auch innerhalb ihrer die Stufen und Grade des Unendlichen unterscheiden kann?
\par

— Wenn der Punkt als Repräsentant des Endlichen gilt, so sind Linien, Flächen, Körper Unendlichkeiten, aber offenbar nicht gleichwerthige Unendlichkeiten, denn die Linie mit ihren unendlich vielen Punkten hat noch unendlich oft in der Fläche, diese unendlich oft im Körper Platz. Um einen formelhaften Ausdruck für diese Unendlichkeiten verschiedener Dimension zu haben, knüpft man an eine höchst einfache Construction an: ein Quadrat von der Seitenlänge \Math{l} hat den Flächeninhalt \Math{l \times l = l^2}, ein Würfel von derselben Kantenlänge \Math{l} hat das Volumen \Math{l \times l \times l = l^3}. Oder, etwas mehr begrifflich ausgedrückt: wenn Linien von der Dimension \Math{l} sind, so sind Flächen von der Dimension \Math{l^2}, Körper von der Dimension \Math{l^3}. Bezeichnen wir demnach die lineare Unendlichkeit mit dem Symbol \Math{\infty}, so ist \Math{\infty^2} das Symbol der flächenhaften, \Math{\infty^3} das der körperlichen Unendlichkeit, oder, indem wir geradezu \Math{\infty} als eine Zahl auffassen, es sind in der Linie \Math{\infty}, in der Fläche \Math{\infty^2}, im Körper \Math{\infty^3} Punkte enthalten. Die Unendlichkeiten verschiedener Dimension sind also zu unterscheiden wie die \Sperr{Potenzen} einer Grösse \Math{\infty}, und hierbei gelten die gewöhnlichen Rechenregeln für Potenzen; so ist beispielsweise \Math{\infty^2 = \infty \times \infty} und \Math{\infty^3 = \infty \times \infty^2 = \infty^2 \times \infty}, woraus man schliessen kann, dass sich die Fläche in \Math{\infty} Linien, der Körper in \Math{\infty} Flächen oder \Math{\infty^2} Linien anordnen lässt. Ein vierdimensionaler Raum wäre natürlich eine Unendlichkeit mit \Math{\infty^4} Punkten, \Math{\infty^3} Linien, \Math{\infty^2} Flächen, \Math{\infty} dreidimensionalen Theilräumen und so fort. Indessen wir müssen den Begriff der Mannigfaltigkeit auch ohne Hülfe der geometrischen Anschauung bilden können. Wenn eine Grösse x, sei es unbegrenzt, sei es zwischen gewissen Grenzen, jeden beliebigen Werth annehmen kann, so nennen wir x eine stetige Veränderliche oder Variable und bezeichnen die unendliche Anzahl ihrer möglichen Werthe mit \Math{\infty}. Nehmen wir eine zweite unabhängige Variable y hinzu und fragen nun nach der Anzahl der \Sperr{Werthpaare}, die das \Sperr{Variablenpaar} x, y annehmen kann, so ist diese Anzahl wiederum unendlich, aber offenbar von der Dimension \Math{\infty^2}; denn da man jeden der \Math{\infty} Werthe von x mit jedem der \Math{\infty} Werthe von y combiniren kann, so ist die Anzahl der Combination \Math{\infty \times \infty = \infty^2}. Wäre dagegen y von x \Sperr{abhängig}, sodass mit dem Werthe von x auch derjenige von y bestimmt wäre, so hätte man nur ebensoviele Werthepaare von x, y, wie man Werthe von x hat, also nur \Math{\infty}. Eine Gruppe von drei unabhängigen Variablen x, y, z kann \Math{\infty^3} Werthgruppen annehmen, und so fort. Im Allgemeinen können wir sagen: wenn eine Mannigfaltigkeit mit n unabhängigen Variablen gegeben und das „Individuum“ innerhalb derselben durch einen bestimmten Werthcomplex dieser n Variablen gekennzeichnet ist, so enthält die Mannigfaltigkeit \Math{\infty^n} Individuen. Um ein paar Beispiele zu nennen, so enthält die Zeit \Math{\infty} Augenblicke, das Spectrum \Math{\infty} homogene Linien; es giebt \Math{\infty} Winde der Richtung nach, \Math{\infty^2} Winde der Richtung und Stärke nach; ferner giebt es \Math{\infty^3} klimatische Zustände, wenn wir einen solchen Zustand durch Luftdruck, Temperatur und Feuchtigkeit definiren. — Nach diesen Ausblicken ins mathematische Gebiet ist die Herleitung der ewigen Wiederkunft aus der angeblichen Erschöpfbarkeit der Atomgruppirungen mit zwei Worten zu widerlegen. Die Zeit hat \Math{\infty} Augenblicke, in jedem Augenblick ist ein „Weltzustand“ unterzubringen — und unter Weltzustand verstehen wir einmal, indem wir die ganz primitive Annahme des demokritischen Atomismus zu Grunde legen, die Lage der materiellen Atome im Raume. Wieviel solche Lagen, solche Weltzustände giebt es? Machen wir die allereinfachste Voraussetzung: nur \Sperr{drei} kugelförmige Atome, die sich beliebig bewegen. Ferner möge es nur auf ihre \Sperr{relative} Lage zu einander ankommen; dann ist jeder „Weltzustand“ durch ein gewisses Atomdreieck charakterisirt, und da ein Dreieck durch drei unabhängige Variable \textit{x}, \textit{y}, z (z. B. die drei Seiten, oder zwei Seiten und einen Winkel) bestimmt ist, so giebt es \Math{\infty^3} solcher Dreiecke. Also \Math{\infty^3} Weltzustände und nur \Math{\infty} Augenblicke: d. h. der mögliche Zeitinhalt hat in der Zeit weder einmal, noch gar unendlich viele Male Platz, sondern es kann von ihm nur ein unendlich kleiner Theil (unter \Math{\infty^3} Weltzuständen nur ein Ausschnitt von \Math{\infty}) zeitlich realisirt werden, oder es wären \Math{\infty^2} Zeiten gleich der unserigen erforderlich, um die Gesammtheit der Atomgruppirungen zu erschöpfen. Nun schreite man aber von diesem \Sperr{ABC} von Materialismus zu verwickelteren Annahmen über die Constitution der Materie fort, man betrachte zehn, hundert, Myriaden Atome und zwar nicht nur als Punkte, sondern als kleine Körperchen von beliebiger irregulärer Gestalt, man nehme die physicalischen Erscheinungen hinzu, Wärme, Electricität, Ätherschwingung, man gehe endlich zur organischen Materie und zum psychischen Leben über und sehe, wie die Dimension, der Unendlichkeits-Exponent des Inbegriffs aller Weltzustände immer unermesslicher anschwillt, während die Zeit durchaus nicht über die Umfänglichkeit ihrer \Math{\infty} Augenblicke hinauszusteigern ist — man vergleiche das armselige Gefäss mit dem überschäumend reichen Inhalt und suche die Wahrscheinlichkeit, dass einer der \Math{\infty^\infty} Weltzustände innerhalb der \Math{\infty^1} Zeitaugenblicke auch nur einmal, zweimal, unendlich viele Male realisirt werde!
\par

— Bei allen diesen Betrachtungen spielt die etwaige \Sperr{Endlichkeit} von Raum, Kraft, Materie gar keine Rolle; wir sahen, dass bereits bei drei Atomen, denen wir noch dazu einen beliebig kleinen Raum als Bewegungsgebiet anweisen können, die Anzahl der Atomgruppirungen \Sperr{un}endlich gross ist. Das liegt einfach an der vorausgesetzten Stetigkeit der Variablen innerhalb ihrer Grenzen, gleichgiltig ob diese Grenzen endliche oder unendliche sind; die kleinste Kugel hat immer \Math{\infty^3} Punkte, die unbegrenzte gerade Linie nur \Math{\infty}. Man müsste geradezu diese Voraussetzung der Stetigkeit aufgeben und den Raum als Fachwerk oder Bienenzellenbau denken; dann allerdings brauchte es nur eine endliche Anzahl von Weltzuständen zu geben, aber mit diesen wäre kein Weltverlauf, kein continuirlicher Fluss herzustellen, und wir hätten überhaupt keine Zeitempfindung. Nietzsches materialistischer Beweis für die Nothwendigkeit der ewigen Wiederkunft darf damit als widerlegt gelten — etwas ganz anderes ist die \Sperr{Denkbarkeit} dieser Hypothese, und wiederum etwas anderes ihr dichterischer, ethischer, speculativer Werth. Die ewige Wiederkunft ist eine gewaltige Conception, ein Mysterium, das schon als Möglichkeit aufregt, erschüttert, ungeheure Folgerungen zulässt. Wir treten diesem „abgründlichen Gedanken“ nicht zu nahe, wenn wir seinen oberflächlichen Beweis verwerfen.
\par

